% !TEX program = pdflatex
\documentclass[10pt,a4paper]{article}
\usepackage[T1]{fontenc}
\usepackage{lmodern}
\usepackage[margin=22mm]{geometry}
\usepackage{amsmath,amssymb,amsthm,booktabs,xcolor}
\usepackage{fancyhdr}
\usepackage[colorlinks=true,linkcolor=teal,urlcolor=teal]{hyperref}
\definecolor{teal}{HTML}{0B8298}
\definecolor{muted}{HTML}{63747A}
\hypersetup{pdftitle={Numerical integration: the trapezoidal rule},pdfauthor={Your Name}}
\pagestyle{fancy}
\fancyhf{}
\fancyhead[L]{\small\color{muted}Math Notes}
\fancyhead[R]{\small\color{muted}Numerical analysis}
\fancyfoot[L]{\small\href{https://example.com}{example.com}}
\fancyfoot[R]{\small\thepage}
\setlength{\headheight}{14pt}
\setlength{\parindent}{0pt}
\setlength{\parskip}{4pt}
\newtheorem{theorem}{Theorem}
\begin{document}
{\small\color{teal}NUMERICAL ANALYSIS}\par
{\LARGE\bfseries Numerical integration: the trapezoidal rule}\par
{\small\color{muted}A short mathematical note}\par
\vspace{6pt}

\section{From a line segment to a quadrature rule}
Suppose we want to approximate
\[
I(f)=\int_a^b f(x)\,\mathrm{d}x.
\]
Replace the graph of \(f\) by the straight line joining its endpoint values. Integrating this linear interpolant gives the trapezoidal approximation
\begin{equation}
T_1(f)=\frac{b-a}{2}\bigl(f(a)+f(b)\bigr).
\label{eq:single}
\end{equation}
The rule is exact for affine functions. A curved graph introduces an error that can be controlled by the second derivative.

\section{The composite rule}
Divide the interval into \(n\) equal subintervals, with
\[
h=\frac{b-a}{n},\qquad x_j=a+jh,\qquad j=0,\ldots,n.
\]
Apply \eqref{eq:single} on each subinterval and add the contributions:
\begin{align}
T_n(f)
&=\sum_{j=0}^{n-1}\frac{h}{2}
  \bigl(f(x_j)+f(x_{j+1})\bigr) \notag\\
&=h\left[\frac{f(a)}2+\sum_{j=1}^{n-1}f(x_j)+\frac{f(b)}2\right].
\label{eq:composite}
\end{align}
Interior values occur twice, so their weights become \(h\). Each endpoint occurs only once and receives weight \(h/2\).

\section{An error estimate}
\begin{theorem}
If \(f\in C^2([a,b])\), there is a point \(\xi\in(a,b)\) such that
\begin{equation}
I(f)-T_n(f)=-\frac{(b-a)h^2}{12}f''(\xi).
\label{eq:error}
\end{equation}
Consequently,
\[
|I(f)-T_n(f)|\le
\frac{(b-a)h^2}{12}\max_{x\in[a,b]}|f''(x)|.
\]
\end{theorem}
\begin{proof}
On \([x_j,x_{j+1}]\), the linear interpolation remainder is
\[
f(x)-p_j(x)=\frac{f''(\eta_x)}{2}(x-x_j)(x-x_{j+1}).
\]
The product has one sign on this interval. The weighted mean-value theorem for integrals therefore gives a local error
\(-h^3f''(\xi_j)/12\). Summing over all intervals and using continuity of \(f''\) yields \eqref{eq:error}.
\end{proof}

\newpage
\section{A worked example}
Take \(f(x)=x^2\) on \([0,1]\). The exact integral is \(I(f)=1/3\). Since \(f''=2\), the error formula gives
\begin{equation}
T_n(f)=\frac13+\frac{1}{6n^2}.
\label{eq:example}
\end{equation}
The approximation is too large, as expected for a convex graph: each chord lies above the curve.

\begin{center}
\begin{tabular}{rrr}
\toprule
\(n\) & \(T_n(f)\) & \(T_n(f)-I(f)\)\\
\midrule
1  & 0.50000000 & 0.16666667\\
2  & 0.37500000 & 0.04166667\\
4  & 0.34375000 & 0.01041667\\
8  & 0.33593750 & 0.00260417\\
16 & 0.33398438 & 0.00065104\\
\bottomrule
\end{tabular}
\end{center}

Doubling the number of subintervals divides the error by four. This is the characteristic behavior of a method with second-order accuracy.

\section{Richardson extrapolation}
For a sufficiently smooth integrand, suppose the leading error has the expansion
\[
T_n(f)=I(f)+Ch^2+O(h^4).
\]
Refining the mesh gives
\[
T_{2n}(f)=I(f)+\frac{C}{4}h^2+O(h^4).
\]
The combination
\begin{equation}
R_n(f)=\frac{4T_{2n}(f)-T_n(f)}{3}
=I(f)+O(h^4)
\label{eq:richardson}
\end{equation}
eliminates the leading error term. For the quadratic example, inserting \eqref{eq:example} into \eqref{eq:richardson} gives \(R_n(f)=1/3\) exactly.

\section{Practical reading of the estimate}
\begin{itemize}
\item The regularity assumption matters. A nonsmooth integrand may not exhibit the predicted convergence rate.
\item A small difference between two approximations is evidence of convergence, not by itself a rigorous error bound.
\item Formula \eqref{eq:error} describes discretization error. Floating-point roundoff is a separate source of error.
\end{itemize}

\vfill
{\small\color{muted}This sample PDF contains selectable text, aligned equations, theorem and proof environments, and a table. It is included with the Math Notes static blog framework.}
\end{document}
